\section*{Capítulo \ref{ch:b2:fragment-orbitals} --- Orbitales de fragmentos y moléculas poliatómicas}

\begin{solution}{exo:b2:fragment-orbitals:1}
$h_1 + h_2$ y $h_1 - h_2$. El plano $xz$, bisector del ángulo y
perpendicular a la molécula, intercambia los dos hidrógenos: $h_1 - h_2$ cambia de
signo y $h_1 + h_2$ no.
\end{solution}

\begin{solution}{exo:b2:fragment-orbitals:2}
Seis orbitales de valencia (el 2s y los tres 2p del oxígeno, y los dos 1s de los hidrógenos), ocho electrones
de valencia y cuatro orbitales llenos.
\end{solution}

\begin{solution}{exo:b2:fragment-orbitals:3}
Los niveles de valencia llenos son un orbital $a_1$ y un conjunto de tres orbitales $t_2$
degenerados: dos energías, dos bandas. La banda $t_2$ arranca electrones de tres
orbitales que contienen seis electrones, frente a un orbital y dos electrones para $a_1$.
\end{solution}

\begin{solution}{exo:b2:fragment-orbitals:4}
El enlace $\pi$ exige que los dos orbitales p de los carbonos sean paralelos, lo que sitúa
los planos de ambos \ce{CH2} en un plano común.
\end{solution}

\begin{solution}{exo:b2:fragment-orbitals:5}
El amoníaco pierde un electrón de su par no enlazante, un orbital mayoritariamente no enlazante del
nitrógeno; el metano, de un \omterm{def:b2:diatomic-mos:bonding}{orbital enlazante} $t_2$. Un electrón no enlazante está más alto,
más cerca del nivel atómico, que uno enlazante: se arranca con más facilidad.
\end{solution}

\begin{solution}{exo:b2:fragment-orbitals:6}
El agua lineal tiene dos orbitales p no enlazantes llenos y degenerados; la flexión rebaja uno de
ellos y eleva ligeramente un nivel enlazante: una ganancia neta con ocho electrones. \ce{BeH2}
solo tiene cuatro, que llenan los dos \omterm{def:b2:diatomic-mos:bonding}{orbitales enlazantes}; el orbital que la flexión
rebajaría está vacío, mientras que los niveles enlazantes pierden solapamiento: la forma lineal es la mejor.
\end{solution}

\begin{solution}{exo:b2:fragment-orbitals:7}
Cada enlace C--H o C--C de un carbono contiguo al centro catiónico puede alinearse con el
orbital p vacío y aporta unos $2\beta^2/\Delta$ de estabilización. \ce{CH3+} no tiene ningún
vecino así; los cationes etilo, isopropilo y terc-butilo tienen uno, dos y tres grupos metilo
junto al centro, cada uno con un enlace alineado: la estabilidad crece en ese
orden.
\end{solution}

\begin{solution}{exo:b2:fragment-orbitals:8}
Para dos niveles iguales el desdoblamiento es $2|\beta|$ y los dos electrones $\pi$ ganan $2|\beta|$,
proporcional a $\cos\theta$. Se reduce a la mitad para $\theta = 60^\circ$ y se anula para
$90^\circ$.
\end{solution}

\begin{solution}{exo:b2:fragment-orbitals:9}
El borano tiene seis electrones de valencia: llenan los tres \omterm{def:b2:diatomic-mos:bonding}{orbitales enlazantes}, y el orbital p
del boro perpendicular al plano queda vacío; no se gana nada al
piramidalizarse. El amoníaco tiene dos más: ocupan el orbital que es no enlazante en
la forma plana y, como en el agua, la piramidalización le permite mezclarse con el orbital 2s y
la combinación de los hidrógenos, y bajar.
\end{solution}

\begin{solution}{exo:b2:fragment-orbitals:10}
Las ecuaciones seculares para $c_1 = c_2 = c_3$ dan $\alpha + 2\beta$; las dos combinaciones
ortogonales a ella dan $\alpha - \beta$ (degeneradas). Dos electrones en $\alpha + 2\beta$: energía
$2\alpha + 4\beta$, más baja que $\ce{H2} + \ce{H+}$ ($2\alpha + 2\beta$), ya que $\beta < 0$.
\end{solution}

\begin{solution}{exo:b2:fragment-orbitals:11}
$\pi_1$: los tres coeficientes del mismo signo, ningún nodo (enlazante); $\pi_2$: signos opuestos
en los carbonos de los extremos y un nodo en el centro (no enlazante); $\pi_3$: signos
alternos, dos nodos (antienlazante). El anión tiene cuatro electrones $\pi$: $\pi_2$ es su
nivel ocupado más alto, con densidad solo en los carbonos de los extremos.
\end{solution}

\begin{solution}{exo:b2:fragment-orbitals:12}
En cada uno de los dos planos que contienen el eje, tres orbitales p paralelos (O, C, O)
dan un \omterm{def:b2:diatomic-mos:bonding}{orbital enlazante}, uno no enlazante (solo sobre los oxígenos) y uno antienlazante: seis
orbitales $\pi$ en total. Los dos enlazantes y los dos no enlazantes están llenos: ocho
electrones $\pi$ (dos enlaces $\pi$ y dos pares no enlazantes del oxígeno en la imagen de Lewis).
\end{solution}

\begin{solution}{pb:b2:fragment-orbitals:1}
\textbf{1.} Ocho: seis del oxígeno y uno de cada hidrógeno.
\textbf{2.} $h_1 + h_2$ y $h_1 - h_2$.
\textbf{3.} 2s y $2\mathrm p_z$.
\textbf{4.} $2\mathrm p_y$.
\textbf{5.} $2\mathrm p_x$, perpendicular al plano molecular.
\textbf{6.} Seis; cuatro llenos.
\textbf{7.} $h_1 - h_2$, que cambia de signo, como $2\mathrm p_z$, al atravesar el plano
perpendicular al eje que pasa por el oxígeno.
\textbf{8.} $h_1 + h_2$.
\textbf{9.} $2\mathrm p_x$ y $2\mathrm p_y$, perpendiculares al eje: una pareja no enlazante
degenerada.
\textbf{10.} (2s, $h_1 + h_2$) enlazante$^2$, ($2\mathrm p_z$, $h_1 - h_2$) enlazante$^2$, y después
$2\mathrm p_x^2\,2\mathrm p_y^2$.
\textbf{11.} No: la pareja degenerada está llena, y todos los electrones están apareados.
\textbf{12.} A la energía de un orbital 2p del oxígeno: es no enlazante.
\textbf{13.} $2\mathrm p_y$ (con el eje de la molécula lineal a lo largo de $z$ y la flexión en $yz$).
\textbf{14.} Se mezcla con 2s y con $h_1 + h_2$ y baja: el nivel $3a_1$.
\textbf{15.} El nivel enlazante ($2\mathrm p_z$, $h_1 - h_2$), que pierde algo de solapamiento cuando los
hidrógenos se apartan del eje: se convierte en $1b_2$.
\textbf{16.} El nivel lleno que baja gana más de lo que pierde el otro: con ocho
electrones la forma angular es la más estable.
\textbf{17.} $104.48^\circ$, entre $90^\circ$ (enlaces formados solo por orbitales p) y $109.5^\circ$:
el orbital 2s participa en parte en el enlace.
\textbf{18.} $2\mathrm p_x$, que sigue siendo el $1b_1$ no enlazante.
\textbf{19.} Cuatro.
\textbf{20.} Del $1b_1$, el orbital $2\mathrm p_x$ no enlazante del oxígeno.
\textbf{21.} Arrancar un electrón no enlazante apenas cambia los enlaces: el ion tiene la
geometría de la molécula y se excita poca energía vibracional. Arrancar un electrón
enlazante cambia las longitudes y los ángulos de enlace, y la banda se reparte entre
muchos niveles vibracionales.
\textbf{22.} \qty{12.62}{eV} frente a \qty{13.62}{eV}: el oxígeno del agua lleva una carga parcial
negativa tomada de los hidrógenos, que eleva sus niveles 2p.
\textbf{23.} No son dos niveles iguales: sus dos combinaciones son los orbitales $3a_1$ y
$1b_1$, de energías distintas. Ambas imágenes describen la misma densidad electrónica
total; solo los orbitales deslocalizados dan las energías de ionización.
\textbf{24.} El agua tiene \textbf{cuatro} bandas de ionización de valencia, la más baja a
$\boldsymbol{\approx \qty{12.6}{eV}}$.
\end{solution}
